\section{The mathematical definition}\label{sec:11-modules:01-definition:1}

Given a ring \(R\) (Chapter 9), a (left) \textbf{\(R\)-module} is an abelian group \((M, +, 0, -(-))\) together with a scalar action \(R \times M \to M\), written \(r \cdot m\), satisfying:

\[
\begin{aligned}
\text{(M1)}&\quad r \cdot (m + n) = r\cdot m + r \cdot n \\
\text{(M2)}&\quad (r + s) \cdot m = r \cdot m + s \cdot m \\
\text{(M3)}&\quad (r \cdot s) \cdot m = r \cdot (s \cdot m) \\
\text{(M4)}&\quad 1 \cdot m = m
\end{aligned}
\]

for all \(r, s \in R\), \(m, n \in M\). This is exactly the vector space definition, with ``field'' replaced by ``ring''. That extra generality is the whole point: \(\mathbb{Z}\)-modules are abelian groups, and \(k[x]\)-modules (for \(k\) a field) are vector spaces equipped with a chosen linear endomorphism. Most importantly for this book, a representation of a quiver \(Q\) is precisely a module over the path algebra \(kQ\); this is why the present chapter is placed immediately before Chapter 12.

\subsection{Sources, quoted}\label{sec:11-modules:01-definition:2}

Formal definitions and citations for this section, gathered here for reference (full entries in the \hyperref[sec:root:bibliography:1]{Bibliography}):

\begin{itemize}
\tightlist
\item
  \textbf{Module.} Given a ring \(R\), a left \(R\)-module is an abelian group under addition together with a scalar action of \(R\) satisfying the left-distributive, right-distributive, and associative-action axioms, plus unitality when \(R\) has an identity (\cite{DummitFoote2003}, Ch. 10 ``Introduction to Module Theory,'' §10.1 ``Basic Definitions and Examples''). This is a structural citation to the section and its numbered axioms, not a verified word-for-word excerpt.
\item
  Weibel (\cite{Weibel1994}) is offered as further reading on modules in the broader context of homological algebra, not an independently verified factual claim.
\end{itemize}
